\documentclass[11pt,a4paper]{article}
\usepackage[T1]{fontenc}
\usepackage[utf8]{inputenc}
\usepackage[scaled=0.95]{helvet}
\renewcommand{\familydefault}{\sfdefault}
\usepackage[margin=1.7cm]{geometry}
\usepackage{enumitem}
\usepackage{titlesec}
\usepackage{xcolor}
\definecolor{accent}{HTML}{6A7BA2}
\definecolor{graytext}{HTML}{6C6C6C}
\definecolor{border}{HTML}{DDDDDD}
\usepackage[colorlinks=true, urlcolor=accent, linkcolor=accent, citecolor=accent]{hyperref}
\usepackage{microtype}
\linespread{0.96}
\pagestyle{empty}
\setlength{\parindent}{0pt}
\titleformat{\section}{\normalfont\large\bfseries}{}{0pt}{}[\vspace{-0.5em}\textcolor{border}{\rule{\textwidth}{0.4pt}}]
\titlespacing*{\section}{0pt}{6pt}{2pt}
\begin{document}
\begin{center}
  {\Huge\bfseries Jan Miksa}\\[3pt]
  {\large Machine Learning Researcher}
  \\[3pt] {\small \href{mailto:jan.miksa@protonmail.com}{jan.miksa@protonmail.com} · +48 669 266 309}
  \\[2pt] {\small\href{https://github.com/OneAndZero24}{GitHub}} · {\small\href{https://www.linkedin.com/in/jan-miksa24/}{LinkedIn}} · {\small\href{https://scholar.google.com/citations?user=bTL7lzgAAAAJ}{Google Scholar}} · {\small\href{https://oneandzero24.github.io/}{Website}}
\end{center}
\vspace{-4pt}
\par\vspace{4pt}\noindent I am a Machine Learning Researcher and PhD Student at GMUM Jagiellonian University, currently working at IDEAS Research Institute. My research focuses on continual learning and machine unlearning, with an emphasis on the geometry and modularity of neural representations. I investigate how the structure of hidden representations can be used to control interference between tasks, preserve learned capabilities, and selectively modify neural network behaviour. My work combines mathematical perspectives on neural representations with empirical studies of deep learning systems, including large-scale diffusion models.
\section*{Selected Publications}
\begin{itemize}[leftmargin=1.2em,topsep=0.5pt,itemsep=1.5pt,parsep=0pt]
\item \textbf{BARRIER: Bounded Activation Regions for Robust Information Erasure}. Jan Miksa, P. Krukowski, P. Spurek, D. D. Rymarczyk, M. Sendera. 2026. ({\color{graytext}Preprint}). [\href{https://arxiv.org/abs/2605.15737}{paper}]
\item \textbf{InTAct: Interval-based Task Activation Consolidation for Continual Learning}. P. Krukowski, Jan Miksa, P. Helm, J. Tabor, P. Spurek. \textit{NeurIPS}. 2026. ({\color{graytext}Accepted}). [\href{https://arxiv.org/abs/2511.17439}{paper}]
\item \textbf{CLASP: Continual Low-rank Adapters for Spatially Placed Concepts from One Hypernetwork}. W. Gromski, P. Krukowski, Jan Miksa, M. Zieba, P. Spurek. 2026. ({\color{graytext}Preprint}). [\href{https://arxiv.org/abs/2610.01331}{paper}]
\item \textbf{The general framework for few-shot learning by kernel HyperNetworks}. M. Sendera, M. Przewięźlikowski, Jan Miksa, M. Rajski, K. Karanowski, et al. \textit{Machine Vision and Applications, 34(4), 53}. 2023. [\href{https://doi.org/10.1007/s00138-023-01403-4}{paper}]
\end{itemize}
\section*{Research Experience}
\noindent\textbf{Senior Machine Learning Research Engineer} \hfill {\small\color{graytext} May 2026 -- Present}\\
\textit{IDEAS Research Institute}\par
\vspace{1pt}\noindent Researching continual learning and machine unlearning, with an emphasis on activation-space methods for controlling interference and preserving model behaviour.\par\vspace{1pt}
\begin{itemize}[leftmargin=1.2em,topsep=0.5pt,itemsep=0.3pt,parsep=0pt]
  \item Developing methods that use the geometry and structure of neural activations to support continual adaptation and targeted information removal.
  \item Investigating activation-level constraints and representation-space interventions in diffusion models and other deep neural networks.
  \item Designing and conducting experiments to evaluate forgetting, retention, and changes in model behaviour.
  \item Building tools for efficient activation tracking and analysis to support research on large neural networks.
\end{itemize}
\vspace{0.5pt}
\noindent\textbf{PhD Researcher} \hfill {\small\color{graytext} July 2025 -- Present}\\
\textit{Jagiellonian University --- GMUM}\par
\vspace{1pt}\noindent Research on continual learning and machine unlearning, focusing on the structure of neural representations and activation-level mechanisms for controlling model updates.\par\vspace{1pt}
\begin{itemize}[leftmargin=1.2em,topsep=0.5pt,itemsep=0.3pt,parsep=0pt]
  \item Investigating representation-space methods for mitigating catastrophic forgetting and enabling targeted changes to learned behaviour.
  \item Developing and evaluating methods based on interval-bounded activations and geometric constraints.
  \item Conducting empirical research on deep neural networks, including diffusion models.
\end{itemize}
\vspace{0.5pt}
\noindent\textbf{Researcher} \hfill {\small\color{graytext} October 2025 -- October 2026}\\
\textit{NCN grant no. 2023/49/B/ST6/01137}\par
\vspace{1pt}\noindent Unlearning in diffusion models.\par\vspace{1pt}
\vspace{0.5pt}
\noindent\textbf{Graduate Research Assistant} \hfill {\small\color{graytext} October 2023 -- July 2025}\\
\textit{Jagiellonian University --- GMUM}\par
\begin{itemize}[leftmargin=1.2em,topsep=0.5pt,itemsep=0.3pt,parsep=0pt]
  \item Conducted research in meta-learning, continual learning, and diffusion models.
  \item Designed neural architectures, custom layers, and end-to-end training pipelines.
  \item Explored hypergradient methods, neural optimizers, and consistency models.
\end{itemize}
\vspace{0.5pt}
\noindent\textbf{Undergraduate Research Assistant} \hfill {\small\color{graytext} July 2022 -- October 2023}\\
\textit{Jagiellonian University --- GMUM}\par
\begin{itemize}[leftmargin=1.2em,topsep=0.5pt,itemsep=0.3pt,parsep=0pt]
  \item Developed Bayesian extensions of meta-learning models.
  \item Designed experiments to investigate uncertainty estimation in learning systems.
\end{itemize}
\vspace{0.5pt}
\section*{Industry Experience}
\noindent\textbf{Computer Vision / Machine Learning Engineer} \hfill {\small\color{graytext} April 2025 -- April 2026}\\
\textit{Aptiv}\par
\begin{itemize}[leftmargin=1.2em,topsep=0.5pt,itemsep=0.3pt,parsep=0pt]
  \item Researched and designed neural network architectures for real-time multi-camera perception.
  \item Implemented custom TensorRT operators and optimized inference performance.
  \item Developed high-performance image-processing pipelines using CUDA and x86 SIMD.
  \item Integrated perception models into real-time systems using low-latency, zero-copy interfaces.
\end{itemize}
\vspace{0.5pt}
\noindent\textbf{Junior Data Scientist} \hfill {\small\color{graytext} September 2024 -- December 2024}\\
\textit{OANDA}\par
\begin{itemize}[leftmargin=1.2em,topsep=0.5pt,itemsep=0.3pt,parsep=0pt]
  \item Developed a Bayesian Marketing Mix Model in PyMC, covering model specification, inference, and validation.
  \item Built an interactive Streamlit interface and prepared deployment on Google Cloud Vertex AI.
\end{itemize}
\vspace{0.5pt}
\noindent\textbf{Data Engineer / Associate Data Engineer} \hfill {\small\color{graytext} March 2023 -- September 2024}\\
\textit{OANDA}\par
\begin{itemize}[leftmargin=1.2em,topsep=0.5pt,itemsep=0.3pt,parsep=0pt]
  \item Designed and maintained transaction-reporting and data-processing pipelines.
  \item Developed a declarative framework for Apache Beam pipelines and optimized cloud resource utilization.
\end{itemize}
\vspace{0.5pt}
\noindent\textbf{Software Developer} \hfill {\small\color{graytext} October 2022 -- February 2023}\\
\textit{EUROLOOP / Willbert}\par
\begin{itemize}[leftmargin=1.2em,topsep=0.5pt,itemsep=0.3pt,parsep=0pt]
  \item Built time-series analytics dashboards (Grafana + InfluxDB) for system monitoring.
  \item Integrated the UI with external systems (payment terminal, invoicing, ads platform).
\end{itemize}
\vspace{0.5pt}
\noindent\textbf{Junior Software Developer} \hfill {\small\color{graytext} March 2022 -- October 2022}\\
\textit{EUROLOOP / Willbert}\par
\begin{itemize}[leftmargin=1.2em,topsep=0.5pt,itemsep=0.3pt,parsep=0pt]
  \item Selected the technology and rebuilt the EV charger UI from scratch with a focus on performance and portability.
  \item Designed the UI system architecture supporting concurrent charging sessions.
\end{itemize}
\vspace{0.5pt}
\section*{Teaching}
\noindent\textbf{Teaching Assistant} \hfill {\small\color{graytext} September 2025 -- Present}\\
\textit{Jagiellonian University}\par
\begin{itemize}[leftmargin=1.2em,topsep=0.5pt,itemsep=0.3pt,parsep=0pt]
  \item Tutoring the artificial intelligence olympiad preparatory class of the \href{https://gmum.net/ai-scientific-club/}{Artificial Intelligence Club} at GMUM, which prepares secondary school students for the \href{https://www.oai.edu.pl/}{Artificial Intelligence Olympiad}.
\end{itemize}
\vspace{0.5pt}
\section*{Education}
\noindent\textbf{Doctor of Philosophy in Technical Computer Science}, \textit{Jagiellonian University} \hfill {\small\color{graytext} 2025 -- Present}\par\vspace{2pt}
\noindent\textbf{Master's Degree in Computational Mathematics}, \textit{Jagiellonian University} \hfill {\small\color{graytext} 2023 -- 2025}\par\vspace{2pt}
\noindent\textbf{Bachelor's Degree in Computational Mathematics}, \textit{Jagiellonian University} \hfill {\small\color{graytext} 2020 -- 2023}\par\vspace{2pt}
\section*{Skills}
\noindent\textbf{Core research:} Machine Learning, Deep Learning, Continual Learning, Representation Learning, Computer Vision, Research\par\vspace{1pt}
\noindent\textbf{Methods and topics:} Machine Unlearning, Transfer Learning, Meta-Learning, Generative Models, Diffusion Models, Neural Networks\par\vspace{1pt}
\noindent\textbf{Research implementation:} PyTorch, Python, CUDA, Model Optimization, Experimental Design, Scientific Computing\par\vspace{1pt}
\noindent\textbf{Systems engineering:} C++, TensorRT, High-Performance Computing, GPU Computing, Computer Architecture, Linux, Docker\par\vspace{1pt}
\noindent\textbf{Additional experience:} Rust, SQL, Google Cloud Platform (GCP), Apache Beam, Data Engineering, Bayesian Statistics\par\vspace{1pt}
\end{document}
